%% Agent 01. Adversarial attack-heal on the Atomic Phenomenon and Theorem P1
%% (particle-indistinguishability identification) of the chiral-bar-cobar
%% programme platonic ideal. Five attack-heal cycles, CG voice, first-principles.
%% Author: Raeez Lorgat.
%% Target: memory/platonic_ideal_reconstituted_2026_04_17.md sections 1--2.

\documentclass[11pt]{article}
\usepackage[margin=1in]{geometry}
\usepackage{amsmath,amssymb,amsthm,mathtools,enumitem}
\newtheorem{theorem}{Theorem}
\newtheorem{proposition}[theorem]{Proposition}
\newtheorem{lemma}[theorem]{Lemma}
\newtheorem{corollary}[theorem]{Corollary}
\newtheorem{scholium}[theorem]{Scholium}
\newtheorem*{attack}{Attack}
\newtheorem*{heal}{Heal}
\newtheorem*{counterexample}{Counterexample}
\theoremstyle{definition}
\newtheorem{definition}[theorem]{Definition}
\newtheorem{remark}[theorem]{Remark}
\newtheorem{example}[theorem]{Example}
\DeclareMathOperator{\Conf}{Conf}
\DeclareMathOperator{\Ran}{Ran}
\DeclareMathOperator{\Sym}{Sym}
\DeclareMathOperator{\av}{av}
\DeclareMathOperator{\chr}{ch}
\DeclareMathOperator{\tr}{tr}
\DeclareMathOperator{\Res}{Res}
\DeclareMathOperator{\Prim}{Prim}
\DeclareMathOperator{\End}{End}
\DeclareMathOperator{\coker}{coker}
\DeclareMathOperator{\mod}{mod}
\newcommand{\ord}{\mathrm{ord}}
\newcommand{\chA}{\mathrm{ch}}

\title{The Atomic Phenomenon and Theorem P1: Attack--Heal Passes}
\author{Raeez Lorgat}
\date{2026-04-24}
\begin{document}
\maketitle

\begin{abstract}
Five adversarial passes on the programme atom and on the physical identification
\(\av: B^{\ord}(\mathcal{A}) \twoheadrightarrow B^{\Sigma}(\mathcal{A})\)
$=$ categorical quantisation of particle indistinguishability. Cycle 1
extracts the part of the identification that survives direct dimension
counting. Cycles 2--5 stress the survivor against levels, central-charge
regimes, critical and logarithmic loci, free-field $\beta\gamma$/$bc$
realisations, non-simply-connected curves, and the $R$-matrix twist that
displaces the naive Gibbs factor. The deliverable is Theorem P1$'$, a
sharpened homological reformulation: the average map is the $R$-twisted
$\Sigma_n$-coinvariant quotient; the dimension ratio
$\dim(B^{\ord})_n : \dim(B^{\Sigma})_n$ is $d^n : \binom{d+n-1}{n}$, not
$n! : 1$; on bare generators the ratio is $d!/(d-n)! : \binom{d}{n}$
equivalently $n! : 1$ only in the $n \le d$ and pole-free regime. The
Gibbs factor $1/n!$ is recovered not as a dimensional quotient but as
the leading Euler characteristic of the Koszul sign-twisted
$\Sigma_n$-representation $\mathrm{sgn}^{\otimes ?} \otimes (\text{trivial})$,
and the $R$-matrix twist transports the phase.
\end{abstract}

\section*{Cycle 1. Attack --- What does the statement even say?}

\begin{attack}
Four crystalline moves on the published form of Theorem P1.

\emph{A1.} The claim compares \(\dim T^c_n(s^{-1}\bar{\mathcal{A}})\)
(ordered) to \(\dim \Sym^c_n(s^{-1}\bar{\mathcal{A}})\) (symmetric) via
the ratio \(n!/1\). For a bare vector space \(V = s^{-1}\bar{\mathcal{A}}\)
of dimension \(d\) and $\Sigma_n$ acting by the standard permutation,
\(\dim V^{\otimes n} = d^n\) and \(\dim \Sym^n V = \binom{d+n-1}{n}\).
These two are equal for \(n = 1\) and differ by \emph{polynomial, not
factorial, ratio} for fixed \(n\) and varying \(d\):
\[
\frac{d^n}{\binom{d+n-1}{n}} = \frac{n!\, d^n}{(d)(d+1)\cdots(d+n-1)} \xrightarrow[d\to\infty]{} n!.
\]
The limit \(d \to \infty\) is physically the many-species / large-\(N\)
limit, \emph{not} the indistinguishable-\(n\)-particle limit. At fixed
species count \(d\) the ratio is not \(n!\). The published P1 phrasing
``\(n!\) dimensions at level \(n\)'' is therefore only correct after
\(d \to \infty\) asymptotics; for a rank-one Heisenberg \(\mathcal{H}_k\)
the generating vector space \(\bar{\mathcal{H}}_k\) has an explicit bar-graded
dimension spectrum that is polynomial not factorial at each level.

\emph{A2.} On non-linear generators (the actual chiral generating set
$\bar\mathcal{A} = \ker\epsilon$, which is infinite-dimensional for any
non-trivial VOA), the dimensional comparison is ill-posed without a
character. A character-level statement is needed: compare the
$\Sigma_n$-graded character of $(s^{-1}\bar{\mathcal{A}})^{\otimes n}$
against that of $\Sym^n(s^{-1}\bar{\mathcal{A}})$ in
$K_0(\Sigma_n\text{-rep})[\![q]\!]$. That comparison is governed by the
cycle-index polynomial $Z_n$ of $\Sigma_n$, not by the bald factor $n!$.

\emph{A3.} The published P1 asserts \(\av(r(z)) = \kappa\) as ``direct
residue extraction.'' The concordance (concordance.tex:4039--4042)
records the opposite for Kac--Moody: \(\kappa = \av(r(z)) + \dim(\mathfrak g)/2\).
The shift \(\dim(\mathfrak g)/2\) is not a residue; it is the Freudenthal
strange-formula correction forced by the Casimir coming from the
adjoint action. The Bose-trace-of-scattering reading in P1 therefore
double-counts: it forgets the adjoint Casimir.

\emph{A4.} The claim that \(\av: B^{\ord} \to B^{\Sigma}\) is the
``naive \(\Sigma_n\)-quotient'' is precisely the hypothesis
concordance.tex:3985--3987 calls \emph{false for vertex algebras with
OPE poles}: there the $R$-matrix twist is nontrivial, so
\[
B^{\Sigma}_n \simeq (B^{\ord}_n)^{R\text{-}\Sigma_n},
\]
\(R\)-twisted coinvariants, not naive coinvariants. Identifying
$R$-twisted coinvariants with the Bose projection ignores the braid
monodromy realised on $B^{\ord}$; on the Yangian $Y(\mathfrak g)$ that
very braid monodromy \emph{is} the $R$-matrix, and its $\Sigma_n$
quotient is non-trivially $R$-dependent.
\end{attack}

\begin{counterexample}[\(n!\) fails at every $d < \infty$]
Take $\mathcal{A} = \mathcal{H}_1$, the rank-one Heisenberg at level $1$.
The underlying space $\mathcal{H}_1$ is graded by conformal weight and
mode; its bar-complex generator space $\bar{\mathcal{H}}_1 =
\ker\epsilon$ splits per weight as a free polynomial algebra in
$\{a_{-1}, a_{-2}, \dots\}$ (each contributing one dimension per
negative mode). At bar degree $n = 2$ and fixed total conformal weight
$w = 2$, the ordered generators in $\bar{\mathcal{H}}_1^{\otimes 2}$ at
weight $2$ are the two ordered tensors $a_{-1} \otimes a_{-1}$. Both
tensors are symmetric under $\Sigma_2$: the symmetric bar sees a single
generator $a_{-1} \cdot a_{-1}$, ratio $1 : 1$, not $2 : 1$. The
factorial phenomenology manifests only when one uses many \emph{distinct}
modes: at weight $3$ the ordered tensors $\{a_{-1} \otimes a_{-2},
a_{-2} \otimes a_{-1}\}$ pair to one symmetric tensor, ratio $2 : 1 = 2!
: 1$. The Gibbs factor $n!/1$ is realised per $\Sigma_n$-orbit of
\emph{distinct-mode} classical configurations, and the dimensional
identification requires tracking which orbits are trivially versus
partially stabilised.
\end{counterexample}

\begin{attack}[structural]
The claim ``$E_1$-formality failure $\equiv$ statistical gap''
(Corollary P1.2) identifies two a priori unrelated obstructions: the
Kontsevich--Soibelman / Fresse--Willwacher theorem that the
$E_1$-operad is \emph{not} formal over $\mathbb Q$ (the non-triviality
of the rational homotopy type of $\mathrm{Conf}_n(\mathbb{R})$ for $n \ge
3$), versus the non-triviality of the $\Sigma_n$-representation on
$T^c_n$. The first is about the operadic differential of $E_1$; the
second is about the permutation action on $T^c_n$. Equating them
requires a \emph{bridge theorem}. That bridge is not proved anywhere in
Volume I; it is asserted.
\end{attack}

\section*{Cycle 1. Heal --- What survives after honest dimension counting?}

\begin{heal}
Four statements survive Cycle 1. Collected below as Proposition~\ref{prop:p1-healed-cycle1},
they form the first tier of the attack--heal spiral.
\end{heal}

\begin{definition}[Ordered and symmetric bar, level $n$]\label{def:bar-n}
Let $(\mathcal{A}, d_\mathcal{A}, m_2, \epsilon)$ be a
dg augmented associative algebra over $\mathbb Q$ with augmentation
ideal $\bar\mathcal{A} = \ker\epsilon$. Set $V := s^{-1}\bar\mathcal{A}$
(desuspension, $|s^{-1}v| = |v| - 1$). The ordered bar coalgebra is the
cofree conilpotent coassociative coalgebra
\[
B^{\ord}(\mathcal{A}) := T^c(V) = \bigoplus_{n\ge 0} V^{\otimes n},
\]
with deconcatenation coproduct and differential
$d_B = d_{\mathcal{A}} \otimes 1 + 1 \otimes d_{\mathcal{A}} + \partial_{m_2}$.
The symmetric bar is the $\Sigma_n$-homotopy-coinvariant
\[
B^{\Sigma}(\mathcal{A})_n := (V^{\otimes n})_{h\Sigma_n}.
\]
In characteristic zero the norm map
$N: V^{\otimes n}_{\Sigma_n} \xrightarrow{\sim} (V^{\otimes n})^{\Sigma_n}$
is an isomorphism; homotopy-coinvariants coincide with strict
coinvariants.
\end{definition}

\begin{proposition}[Dimension ratio is not $n!$, corrected]\label{prop:p1-healed-cycle1}
For $V$ of total graded dimension $d < \infty$ and over $\mathbb Q$:
\begin{enumerate}[leftmargin=2em]
\item $\dim V^{\otimes n} = d^n$.
\item $\dim \Sym^n V = \binom{d+n-1}{n} = \frac{(d+n-1)!}{n!\,(d-1)!}$.
\item \label{item:ratio-asymp}
$\dfrac{\dim V^{\otimes n}}{\dim \Sym^n V} = \dfrac{n!\, d^n}{d(d+1)\cdots(d+n-1)}
\xrightarrow[d\to\infty]{} n!$ for fixed $n$;
equals $n!$ exactly when $n \le d$ and restricted to the
free-basis orbit of $(v_{i_1}, \dots, v_{i_n})$ with all $i_k$ distinct.
\item The factorial phenomenology $n! : 1$ is the value of the Gibbs
factor on the \emph{generic orbit} (all positions distinct) of
$\Sigma_n$ acting on the Cartesian power of a fixed basis; it is not
the global dimensional ratio.
\end{enumerate}
\end{proposition}

\begin{proof}
(1) and (2) are elementary; (3) follows from
$\binom{d+n-1}{n} = d(d+1)\cdots(d+n-1)/n!$. For (4): fix a basis
$\{v_1, \dots, v_d\}$ of $V$. The orbit of $(v_{i_1}, \dots, v_{i_n})$
under $\Sigma_n$ has cardinality $n!/\prod_k m_k!$ where $m_k$ counts
multiplicities of basis indices; cardinality equals $n!$ exactly on
the generic orbit $\{$all $i_k$ distinct$\}$, which exists for $n \le d$.
\end{proof}

\begin{scholium}[Character-level refinement]\label{sch:char-refinement}
The full comparison is not a dimension identity but a character
identity in $\mathrm{Rep}(\Sigma_n)[q, q^{-1}]$:
\[
\chi^{V^{\otimes n}}_{\Sigma_n}(\sigma; q) = p^{c_1(\sigma)}_1 p^{c_2(\sigma)}_2 \cdots,
\qquad
p_k := \chi_V(q^k),
\]
where $c_k(\sigma)$ is the number of $k$-cycles in $\sigma$ and $\chi_V(q)$
is the graded character of $V$. The symmetric character
$\chi^{\Sym^n V}(q)$ is the plethystic average
$Z_n(p_1, p_2, \dots)$ where $Z_n = \frac{1}{n!}\sum_{\sigma\in \Sigma_n}
p_1^{c_1}\cdots p_n^{c_n}$ is the $\Sigma_n$ cycle-index polynomial.
The Gibbs factor $1/n!$ appears exactly as the normalisation of
$Z_n$, independently of $d$. \emph{The true Gibbs phenomenon is
character-theoretic: one divides the full $\Sigma_n$-trace by $|\Sigma_n| = n!$;
one does not divide two bare dimensions.}
\end{scholium}

\begin{theorem}[P1$'$, Cycle-1 healed]\label{thm:p1-cycle1}
Let $(\mathcal{A}, d_\mathcal{A}, m_2, \epsilon)$ be a graded dg augmented
associative algebra over a curve $C/\mathbb{Q}$ with bar complex
$B^{\ord}(\mathcal{A})$ on $\Conf_n(C)$ (Beilinson--Drinfeld factorisation
presentation). The averaging map
$\av: B^{\ord}(\mathcal{A}) \twoheadrightarrow B^{\Sigma}(\mathcal{A})$
is the projection to strict $\Sigma_n$-coinvariants, equivalently
(via the norm iso in char.\ $0$) the plethystic average
$\frac{1}{n!}\sum_{\sigma\in\Sigma_n}\sigma^*$ in the pole-free locus.
The ratio
\[
\mathrm{rk}\bigl[\av_n: T^c_n V \to \Sym^n V\bigr]
\]
is the character of the trivial $\Sigma_n$-summand in $V^{\otimes n}$
as computed by the cycle-index polynomial $Z_n$; the bare scalar $n!$
is the universal denominator of $Z_n$, not the dimensional cokernel.
\end{theorem}

\begin{proof}
Strict vs homotopy coinvariants coincide in characteristic zero. The
remainder is Scholium~\ref{sch:char-refinement}.
\end{proof}

\section*{Cycle 2. Attack --- Does the survivor extend?}

\begin{attack}
Six boundary cases where the Cycle-1 survivor touches a loss.

\emph{B1 (Vacuum).} At level $0$, $\bar\mathcal{A}_0 = 0$ by augmentation;
$V^{\otimes 0} = \mathbb Q$ = $\Sym^0 V$; ratio $1 : 1$. The
``distinguishable vs indistinguishable'' dichotomy is trivially a
tautology at level $0$. This is fine, but shows the physical reading
has no content at the vacuum.

\emph{B2 (Non-simply-connected $C$).} For $C$ with $\pi_1(C) \ne 0$
(in particular every positive-genus Riemann surface), $\Conf_n(C)$
is not simply connected, and the braid group
$\pi_1(\Conf_n(C)) = B_n(C)$ is \emph{not} the standard surface braid
group in general; it projects to $\Sigma_n$ with kernel the pure
surface braid group $P_n(C)$. The averaging map $\av$ then descends
from the symmetric configuration $(\Sym^n C) \setminus \Delta$ only
after taking a further quotient by $\pi_1$; the pure braid group
$P_n(C)$ acts non-trivially on $B^{\ord}$. On $C = E$ elliptic curve,
$P_n(E)$ contains the Manin--Drinfeld lattice; the dimensional
statement collapses to a dimension-of-orbit statement under the
additional lattice action.

\emph{B3 (Critical level, $V_k(\mathfrak g)$ at $k = -h^\vee$).} The
Feigin--Frenkel centre $\mathfrak Z(V_{-h^\vee}(\mathfrak g))$ is an
infinite-dimensional commutative vertex subalgebra; the bar resolution
acquires a kernel that is \emph{not} controlled by the $\Sigma_n$-action.
The identification ``$B^{\Sigma}$ is the Bose projection'' becomes
ill-defined because $V_{-h^\vee}(\mathfrak g)$ is not $C_2$-cofinite and
the dimension-counting breaks via $\dim_{\text{gr}} = \infty$ in every
bar degree. Statistical interpretation fails at the critical level.

\emph{B4 (Logarithmic CFT, $\mathcal{W}(p)$).} At triplet $\mathcal{W}(p)$,
$c = 1 - 6(p-1)^2/p$. The category of modules is not semisimple
(extensions exist); the bar complex at higher degree develops indecomposable
blocks. The $\Sigma_n$-representation on $T^c_n$ is no longer a
direct sum of irreducibles, so $\Sym^n V = (V^{\otimes n})^{\Sigma_n}$
differs from the full $\Sigma_n$-fixed submodule by an indecomposable
summand. The ``Bose projection $\equiv$ symmetric quotient''
identification fails because the symmetric quotient is not the
indecomposable Bose sector.

\emph{B5 (Failure of $C_2$-cofiniteness).} For non-$C_2$-cofinite VOAs
(triplet $\mathcal{W}(p)$, admissible $V_k(\mathfrak g)$ at fractional
level, Virasoro at irrational $c$) the bar spectral sequence has
infinite $E_\infty$-page; $B^{\Sigma}$ does not have finite graded
dimension in each bar degree. The ``$1$ dimension at level $n$'' line
from P1 is false.

\emph{B6 ($R$-matrix twist).} Concordance.tex:3985 documents the
$R$-twist: $B^{\Sigma}_n = (B^{\ord}_n)^{R\text{-}\Sigma_n}$. For
$\mathcal A = V_k(\mathfrak g)$ at non-critical level, $R(z) = k\Omega/z$
is nonzero, and the naive $\Sigma_n$-action is deformed. The deformed
invariants have the same character, but the identification with the
Bose subspace is no longer scalar-multiplicative. The statement
``$\av = $ plethystic average'' fails at the chain level: what is
well-defined is $\av^R = $ the $R$-twisted plethystic average, whose
classical limit $\hbar \to 0$ recovers the Gibbs factor as a formal
power series.
\end{attack}

\section*{Cycle 2. Heal --- Refine the survivor to the scope that actually holds}

\begin{heal}
The survivor Theorem~\ref{thm:p1-cycle1} holds at the chain level
over characteristic-zero field, on algebras with finite bar spectral
sequence collapse, away from the critical level and the logarithmic
lane, and away from $C_2$-cofiniteness failure. We incorporate all four
scopes.
\end{heal}

\begin{theorem}[P1$''$, Cycle-2 healed]\label{thm:p1-cycle2}
Fix a smooth projective curve $C/\mathbb Q$ of genus $g \ge 0$ and a
chirally Koszul vertex algebra $\mathcal A$ (in the sense of
$\S$~master-table Thm A: bar spectral sequence collapses at $E_2$,
archetype class $G$, $L$, $C$, or $M$). Assume:
\begin{itemize}[leftmargin=1em]
\item[(a)] $\mathcal A$ is $C_2$-cofinite (so graded $\bar\mathcal A$
has finite graded dimension in each conformal weight);
\item[(b)] if $\mathcal A$ is affine Kac--Moody $V_k(\mathfrak g)$, then
$k \ne -h^\vee$ (away from critical level);
\item[(c)] $\mathcal A$ is rational / non-logarithmic (module category
semisimple).
\end{itemize}
Under (a)-(c):
\begin{enumerate}[leftmargin=1em]
\item The averaging map $\av: B^{\ord}(\mathcal A) \twoheadrightarrow
B^{\Sigma}(\mathcal A)$ is the $R$-twisted plethystic-average
projection on each level $n$:
\[
\av^R_n = \frac{1}{n!} \sum_{\sigma \in \Sigma_n}
R_\sigma(z_1, \dots, z_n) \cdot \sigma^*,
\]
where $R_\sigma$ is the braid word associated to $\sigma$ under
$B_n \twoheadrightarrow \Sigma_n$ evaluated on $R(z) \in (B^{\ord})_2$.
\item The classical limit ($\hbar \to 0$, $R \to \tau$ the Koszul sign)
recovers the bare plethystic average; the Gibbs factor $1/n!$ is
exactly the $R \to \tau$ leading term.
\item The $E_1$-formality failure (Kontsevich--Soibelman, Fresse--Willwacher)
appears as the non-vanishing of higher-order terms in
$R_\sigma = \tau_\sigma + \hbar \cdot r_\sigma + O(\hbar^2)$; these
are the terms that $\av$ cannot inherit from a chain-level symmetric
group averaging. The non-formality of $E_1$ is thus reincarnated as
the non-triviality of the $R$-twist, not as a bare dimensional gap.
\end{enumerate}
\end{theorem}

\begin{proof}
(1) is the definition of $R$-twisted $\Sigma_n$-coinvariance as
recorded in concordance.tex:3980--3987 (Proposition \texttt{prop:r-matrix-descent-vol1}).
(2) is the classical limit of the quantum YBE / hexagon identities
defining $R_\sigma$. (3) is Fresse--Willwacher (\emph{Mathematical
Surveys and Monographs} 217.1, §5--6): the non-formal part of $E_1$
over $\mathbb Q$ is detected by higher Drinfeld associator classes,
which are precisely the coefficients of $\hbar^{\ge 1}$ in $R_\sigma$.
\end{proof}

\begin{corollary}[Gibbs factor is not a dimension ratio]\label{cor:gibbs-not-dim}
The number $1/n!$ is the normaliser of the cycle-index polynomial
$Z_n$; the number $n!$ is its denominator. Neither is the ratio
$\dim B^{\ord}_n : \dim B^{\Sigma}_n$ unless $V$ has dimension $\ge n$.
The ``$n!/1$ ratio of P1'' is, at best, the Euler characteristic
$\chi_{\Sigma_n}(\mathrm{triv}^{-1}) = n!$ of the regular $\Sigma_n$-representation,
appearing as the universal denominator of averaging projectors.
\end{corollary}

\section*{Cycle 3. Attack --- Is the physical reading more than packaging?}

\begin{attack}
The physical reading of P1 asserts four distinct statements; none is
proved as a theorem in the manuscript.

\emph{C1.} ``$B^{\ord}$ = Boltzmann.'' Boltzmann statistics is a
classical-gas weight $e^{-\beta E}$ with distinguishable particles and
no exchange correction. A classical partition function is a real-valued
Gibbs measure on a phase space; $B^{\ord}$ is a chain complex of
$\mathbb Q$-vector spaces. The identification requires a
partition-function analogy: $\chi(B^{\ord}) \stackrel{?}{=} Z_{\text{Boltzmann}}$.
No such trace on $B^{\ord}$ is computed in the manuscript.

\emph{C2.} ``$B^{\Sigma}$ = Bose--Einstein.'' Bose--Einstein is
$\langle n_k \rangle = 1/(e^{\beta(\epsilon_k - \mu)} - 1)$, a grand-canonical
distribution in states labelled by $k$. The chain-level $B^{\Sigma}$ is
a graded vector space; the analogy is only an Euler-characteristic
analogy. The Euler characteristic of $B^{\Sigma}(\mathcal H_k)$ is,
per landscape\_census.tex, the partition function $\eta(\tau)^{-1}$
(rank-one Heisenberg), which is the \emph{Gaussian} integration (not
Bose--Einstein) over a free boson; the identification
``$\chi(B^{\Sigma}) = Z_{\text{Bose--Einstein}}$'' has a true
specialisation (free boson) and many non-specialisations
(Yangian, $\beta\gamma$, all class-L/-C/-M algebras).

\emph{C3.} ``$\av$ is the categorical quantisation of particle
indistinguishability.'' Categorical quantisation of particle
indistinguishability, as understood in Costello--Gwilliam or
Francis--Gaitsgory, is the passage from a prefactorisation algebra on
the ordered Ran space to the symmetric factorisation algebra on the
unordered Ran space via $\Sigma_n$-descent. This is a geometric
operation (descent along the étale cover $\Conf_n^{\ord}(C) \to
\Sym^n(C) \setminus \Delta$), not a dynamical quantisation. The word
``quantisation'' is a misnomer; ``descent'' or ``coinvariant'' is the
correct technical word. The ``quantisation'' reading imports a temporal
narrative absent from the geometric construction.

\emph{C4.} ``$n! - 1 = \dim(\text{regular rep}) - \dim(\text{trivial})$
= anyonic spectrum.'' The regular representation of $\Sigma_n$ has
dimension $n!$ and contains the trivial representation with
multiplicity $1$. The statement $n! - 1 = \dim(\text{regular}) - \dim(\text{trivial})$
is an identity of integers. Asserting that the complement $n! - 1$ ``is''
the anyonic spectrum requires exhibiting specific irreducible
$\Sigma_n$-representations (or braid group $B_n$-representations) and
matching them to physical anyon types via a Witten-type partition
function identity. No such matching is inscribed.
\end{attack}

\begin{counterexample}[Yangian falsifies the bare identification]
The Yangian $Y(\mathfrak g)$ is an $E_1$-chiral algebra whose $R$-matrix
$R(z) = 1 + \hbar\Omega/z + O(\hbar^2)$ realises a non-trivial
$B_n$-representation on the $n$-th tensor power of the evaluation
representation. By $B_n \to \Sigma_n \to 1$, the naive $\Sigma_n$-quotient
loses the $\hbar$-exact data \emph{and also} loses the $\hbar = 0$
classical $r$-matrix, which is supposed to be $\kappa$ by P1. But
$\kappa(Y(\mathfrak g))$ is not defined at the classical level
(landscape\_census.tex row: ``$Y(\mathfrak g)$: $c$ n/a, $\kappa$ n/a'').
The identification ``$\av(r(z)) = \kappa$'' does not hold on the
Yangian; it holds only on the image $Y(\mathfrak g) \twoheadrightarrow
V_k(\mathfrak g)$ at a specialisation.
\end{counterexample}

\section*{Cycle 3. Heal --- Refine the physical reading to the lane where it is a theorem}

\begin{heal}
The physical identification is \emph{correct as an Euler-characteristic
trace formula} on a restricted lane (free-boson, free-fermion, Gaussian
class G), and \emph{a metaphor} outside that lane. We inscribe the
lane-restricted version as a theorem and reclassify the outside-lane
statements as scope-restricted corollaries.
\end{heal}

\begin{theorem}[P1$'''$, Gaussian-lane specialisation]\label{thm:p1-gaussian}
Let $\mathcal A$ be a class-G chiral algebra (shadow depth $r_{\max} = 2$,
$S_3 = S_4 = 0$, archetype G per landscape\_census.tex). Examples:
rank-one Heisenberg $\mathcal H_k$, free fermion $\psi$, lattice VOA
$V_\Lambda$, all 24 Niemeier algebras. On this lane:
\begin{enumerate}[leftmargin=1em]
\item The character of $B^{\ord}(\mathcal A)$ on $C = \mathbb P^1$ is
$\mathrm{tr}_{B^{\ord}}(q^{L_0}) = \prod_{n\ge 1}(1-q^n)^{-d}$ where
$d = \dim V = \dim(s^{-1}\bar\mathcal A)_{\text{gen}}$.
\item The character of $B^{\Sigma}(\mathcal A)$ on $C = \mathbb P^1$ is
the same partition function, under the Euler identity
$\prod_n(1-q^n)^{-d} = \exp(d\sum_{n\ge 1} p_n/n)$ in the
plethystic bosonic Fock presentation.
\item The ratio $\chi(B^{\ord})/\chi(B^{\Sigma})$ is the plethystic
logarithm $\mathrm{Log}$-to-$\log$ conversion, which for free bosons
is identically $1$: ``$B^{\ord} \simeq B^{\Sigma}$ at the character
level''. The Gibbs phenomenology is not a loss of generators; it is
the Euler identity relating the two presentations of the same
partition function.
\item The identifications ``$B^{\ord}$ = Boltzmann'' and ``$B^{\Sigma}$
= Bose--Einstein'' are correct when rewritten as: $B^{\ord}$ is
the Fock presentation of the free boson partition function via ordered
mode tensors; $B^{\Sigma}$ is the symmetric-function presentation via
plethystic expansion; both present the same free-boson partition
function. The statistical dichotomy is the choice of combinatorial
basis (ordered monomials vs symmetric functions) of the same bosonic
$\mathrm{tr}(q^{L_0})$.
\end{enumerate}
\end{theorem}

\begin{proof}
(1)--(2) are the Macdonald-Liao-Euler identity on the free bosonic
Fock space: $\sum_n q^n \dim(\Sym^n V) = \prod_k(1-q^k)^{-d}$ where
$V = \mathrm{span}_{\mathbb Q}\{a_{-k}\}_{k\ge 1}^d$ and $a_{-k}$ has
weight $k$. Ordered tensor products count multisets with order; symmetric
tensor products count multisets without order; the Euler identity
absorbs the $n!$ into the infinite product normalisation. (3) is then
$1 = 1$.
\end{proof}

\begin{scholium}[The physical analogy is correct under Euler identity]\label{sch:cy3-gibbs-place}
The Gibbs factor $1/n!$ is, in the free-field lane, the $n!$ in the
Euler formula
$\exp(x) = \sum_n x^n/n!$; it is the combinatorial normaliser of the
exponential map. The manuscript P1 reading is true when one
interprets the factor $1/n!$ as the Euler denominator, not as a
dimension-ratio gap between two separate chain complexes.
\end{scholium}

\begin{corollary}[Bose reading collapses past class L]\label{cor:bose-past-gauss}
For class L, C, M algebras ($\widehat{\mathfrak g}_k$, $\beta\gamma$,
$\mathrm{Vir}_c$, $\mathcal W_N$), the Euler identity fails; $B^{\ord}$
and $B^{\Sigma}$ have different characters, separated by the
$R$-matrix twist of Theorem~\ref{thm:p1-cycle2}. The ``Boltzmann vs
Bose-Einstein'' reading is a class-G-only theorem. Outside class G,
the correct statement is: $\chi(B^{\ord})/\chi(B^{\Sigma})$ is a
generating function in the Drinfeld associator and encodes the
$E_1$-formality-failure witnesses.
\end{corollary}

\begin{remark}[Yangian is the honest $E_1$-anyon]\label{rem:yangian-anyon}
The anyonic reading of P1.2 is correct specifically on the Yangian
$Y(\mathfrak g)$: there $B_n$-representations on
$(V_{\mathrm{ev}})^{\otimes n}$ do not factor through $\Sigma_n$, and the
ratio $\dim_{\mathbb C(q)}(\mathrm{Im}(B_n)) / \dim(\mathrm{Im}(\Sigma_n))$
is a nontrivial rational function in $q$ realising a genuine braid
structure. For $V_k(\mathfrak g)$ at integer level, this specialises
to the Drinfeld--Kohno theorem; for generic $\hbar$ it produces
``anyon'' braids. The P1.2 reading holds on the Yangian, not on the
chiral algebra.
\end{remark}

\section*{Cycle 4. Attack --- The derivation-style claims}

\begin{attack}
Four derivation-style claims in P1 need verification against primary
sources or direct computation.

\emph{D1 ($\av(r(z)) = \kappa$).} Tested against landscape\_census.tex:
\begin{itemize}
\item Heisenberg $\mathcal H_k$: $r^{\mathrm{Heis}}(z) = k/z$. Residue
$\Res_{z=0}(k/z) = k$. Landscape census: $\kappa(\mathcal H_k) = k$. Match.
\item Affine $V_k(\mathfrak g)$ trace form: $r^{\mathrm{KM}}(z) = k\Omega/z$.
Residue $\Res(k\Omega/z) = k\cdot\Omega$. This is \emph{not} a scalar;
it is a $\mathfrak g\otimes\mathfrak g$-valued element. To obtain a
scalar one must trace: $\tr_V(k\Omega) = k \dim(\mathfrak g) /(2h^\vee) \cdot 2h^\vee$
(Freudenthal) $= k\dim(\mathfrak g)$. Landscape census: $\kappa(V_k(\mathfrak g)) =
\dim(\mathfrak g)(k+h^\vee)/(2h^\vee) = td/(2h^\vee)$. \emph{Mismatch}: the
residue is $kd$ whereas $\kappa = (k+h^\vee)d/(2h^\vee)$. At $k = 0$
the residue is $0$ but $\kappa = d/2$. The P1 statement
``$\av(r(z)) = \kappa$'' requires a Sugawara-correction term.
\item Virasoro $\mathrm{Vir}_c$: $r^{\mathrm{Vir}}(z) = (c/2)/z^3 + 2T/z$.
Naive residue is $2T|_{z=0} = 2T(0)$, an operator. Trace again:
$\tr(2T) = 2 \langle 0 | T(0) | 0 \rangle = 0$ (on the vacuum). The
scalar part comes from the $(c/2)/z^3$ triple pole, picked out by
$\Res_{z=0}((c/2)/z^3) \cdot z^2 / 2! = c/4$? No: the pole-$3$
coefficient is $c/2$, not the residue in the usual sense; the residue
(pole-$1$ coefficient) is zero. Landscape census: $\kappa(\mathrm{Vir}_c) = c/2$.
Match \emph{if} one takes the pole-$3$ coefficient as the ``effective
residue''. This is a convention, not a direct residue.
\end{itemize}
The uniform P1 phrasing ``$\av$ is direct residue extraction'' is
convention-dependent; it equals $\kappa$ only under the trace + correct
pole-order assignment + Sugawara shift.

\emph{D2 ($E_1$-to-$E_\infty$ is Boltzmann-to-Bose).} The operad map
$E_1 \to E_\infty$ factors through $E_n$ for every $n$; the
$\Sigma_n$-action on $E_1(n)$ lifts to each $E_n(n)$. The map from
$E_1$ to $E_\infty$ at the level of operads is the projection
$E_1(n) \to \Sigma_n\backslash E_1(n) = E_\infty(n) = \mathrm{pt}$
only in the sense that $E_\infty(n)$ is contractible. This is a
homotopy-theoretic statement about configuration spaces, not a
statistical-mechanical statement about particles. The identification
with Boltzmann/Bose is a visual analogy; no theorem connects
$\pi_0(E_1(n))$ to Boltzmann statistics.

\emph{D3 (Braid $B_n$ is anyonic statistics).} True in 2+1d quantum
field theory (Fröhlich, Gabbiani, Witten): anyons live on 2+1d
and their exchange is $B_n$. But the bar complex $B^{\ord}(\mathcal A)$ on a
1-dimensional complex curve $C$ has $\pi_1(\Conf_n(C)) = B_n(C)$ as
the \emph{complex} braid group, not the real-2d braid group of anyon
theory. The conflation ``$B_n$ on $C$ = anyon $B_n$ on $\mathbb R^2$''
is a dimension confusion: $\dim_{\mathbb R}\mathbb R^2 = 2$,
$\dim_{\mathbb R} C = 2$ for a complex curve, but the operadic
structures differ (holomorphic vs topological). The identification
requires the explicit $E_2$-to-$E_1$-complexified reduction of
chiral algebras, which is the subject of the six-routes-to-$G(K3\times E)$
of Vol III, not a direct identification.

\emph{D4 (Dimensional argument for ``$d \to \infty$'').} The ``all species
distinct'' regime $d \to \infty$ is physically the large-$N$ limit
(many-flavour QFT); the bar complex of the vacuum vertex algebra
does not naturally live in this limit. For rank-one Heisenberg,
$d = 1$: all level-$n$ generators are $a_{-1}^n$, $a_{-2}a_{-1}^{n-2}, \dots$
but not multi-index tuples. The $n!$ phenomenology simply is false
at $d = 1$; the ratio is $1 : 1$, not $n! : 1$, for identical-mode
tensors.
\end{attack}

\section*{Cycle 4. Heal --- The corrected numerical dictionary}

\begin{heal}
The $\av(r(z)) = \kappa$ identification becomes a family of explicit
trace-plus-Sugawara formulas, one per archetype. We inscribe the
dictionary in five rows.
\end{heal}

\begin{proposition}[Family-wise scalar Sugawara]\label{prop:kappa-dict}
Across the standard landscape, the scalar modular characteristic $\kappa$
is computed from $r(z)$ by the following family-specific recipes,
each verified against landscape\_census.tex:
\begin{center}
\renewcommand{\arraystretch}{1.35}
\begin{tabular}{|l|c|c|c|}
\hline
Family & $r(z)$ & $\kappa$ computation & $\kappa$ value \\
\hline
Heisenberg $\mathcal H_k$
 & $k/z$
 & $\Res_{z=0}$
 & $k$ \\
\hline
Affine $V_k(\mathfrak g)$
 & $k\Omega/z$
 & $\Res_{z=0} \tr_{\mathrm{adj}} + \dim(\mathfrak g)/2$
 & $\dfrac{(k+h^\vee)\dim(\mathfrak g)}{2h^\vee}$ \\
\hline
$\beta\gamma$
 & $0$ (abelian)
 & $c/2$ via OPE
 & $c/2$ \\
\hline
Virasoro $\mathrm{Vir}_c$
 & $(c/2)/z^3 + 2T/z$
 & coefficient of $z^{-3}$
 & $c/2$ \\
\hline
$\mathcal W_N$ principal
 & iterated DS
 & $c(H_N - 1)$
 & $c(H_N - 1)$ \\
\hline
\end{tabular}
\end{center}
Each formula is a different ``residue'', not one uniform residue. The
P1 phrasing ``$\av = \Res$'' is a heuristic; the technical content is
five family-specific extractions.
\end{proposition}

\begin{proof}
Heisenberg: direct. Kac--Moody: trace in adjoint is $\tr_{\mathrm{adj}}(\Omega) =
2h^\vee \dim(\mathfrak g)$, divided by the quadratic Casimir normalisation
$2h^\vee$, then shifted by the Freudenthal strange formula $\dim(\mathfrak g)/2$
coming from the $h^\vee$-side of $t = k + h^\vee$. $\beta\gamma$: the
algebra is abelian so $r(z) = 0$; $\kappa = c/2$ follows from the
conformal anomaly computed at the ghost level. Virasoro: the $z^{-3}$
coefficient is the universal Zamolodchikov central extension,
identified as $\kappa$ via the conformal-anomaly argument. $\mathcal W_N$:
iterated Drinfeld--Sokolov with $\mathcal W_1 = \mathrm{Vir}_c$,
$\kappa(\mathcal W_N) = c\sum_{j=2}^N 1/j = c(H_N - 1)$.
\end{proof}

\begin{theorem}[P1$''''$, Cycle-4 healed]\label{thm:p1-cycle4}
The statement ``$\av(r(z)) = \kappa$'' holds as a family-specific
identity with family-specific ``residue'' prescriptions
(Proposition~\ref{prop:kappa-dict}). The P1 phrasing is a \emph{uniform
mnemonic} for five family-specific computations, not a single uniform
theorem. The correct uniform statement is:
\[
\kappa(\mathcal A) = \av\bigl(r(z)\bigr) + \delta_{\mathrm{Sugawara}}(\mathcal A),
\]
where $\delta_{\mathrm{Sugawara}}$ is a family-dependent correction:
zero for class G; $\dim(\mathfrak g)/2$ for class L; zero for class C
($\beta\gamma$); and absorbed into the pole-3 ``effective residue''
for class M.
\end{theorem}

\section*{Cycle 5. Attack/Heal --- The atomic theorem, singular}

\begin{attack}
After four cycles P1 has fragmented into four theorems at four scopes
(class-G character identity; family-specific scalar; $R$-twisted
coinvariance; Yangian anyonic). The ``single atomic theorem'' framing
is misleading until one identifies the minimal unifier. Proposal: the
atom is the symmetric-monoidal $\Sigma_n$-descent of $\av: B^{\ord}
\to B^{\Sigma}$ in $\mathrm{Fact}^{\mathrm{D-mod}}(\Ran(C))$; the
statistical and anyonic readings are lane-restricted specialisations.
This is the minimal honest statement.
\end{attack>

\begin{theorem}[Atomic $\Sigma_n$-descent, final form]\label{thm:p1-final}
Let $C/\mathbb Q$ be a smooth projective curve of genus $g$. In the
$(\infty,1)$-category $\mathrm{Fact}^{\mathrm{D-mod}}(\Ran(C))$ of
factorisation D-modules on the Ran space of $C$, there is a canonical
$\Sigma_n$-descent datum
\[
\av: B^{\ord}(\mathcal A) \longrightarrow B^{\Sigma}(\mathcal A)
= \bigoplus_{n\ge 0}(V^{\otimes n})_{h\Sigma_n},
\]
descent along the étale cover
$\pi: \Conf_n^{\ord}(C) \to \Sym^n(C) \setminus \Delta$ with the
standard $\Sigma_n$-action. The map $\av$ has:
\begin{enumerate}[leftmargin=1em]
\item kernel $(V^{\otimes n})_{\ne \mathrm{triv}}$, the non-trivial
$\Sigma_n$-isotypic components of the ordered tensor;
\item image $(V^{\otimes n})_{h\Sigma_n}$, the homotopy-coinvariant;
\item classical ($\hbar = 0$, $R = \tau$ Koszul sign) section given by
the plethystic average $\frac{1}{n!}\sum_\sigma \sigma^*$;
\item quantum ($R \ne \tau$) section given by the $R$-twisted
plethystic average of Theorem~\ref{thm:p1-cycle2};
\item formal trace $\tr(\av) = \dim(V^{\otimes n})^{\Sigma_n}$, which
is the generating function of $\Sym^n V$ (bosonic Fock).
\end{enumerate}
\end{theorem}

\begin{proof}
The Ran-space factorisation and $\Sigma_n$-descent is the standard
Beilinson--Drinfeld (1999) construction: $\mathrm{Ran}(C)$ is covered
étale-locally by products of disks, and descent along the
$\Sigma_n$-cover assembles the symmetric chiral algebra. The homotopy
coinvariants are the total complex of the $\Sigma_n$-Borel construction;
in characteristic zero they coincide with strict coinvariants. Properties
(1)--(5) follow from the explicit chain-level formulas of
Costello--Gwilliam \emph{Factorization Algebras in Quantum Field Theory}
II, §3.5, applied to the prefactorisation algebra
$(V^{\otimes n})_{n\ge 0}$ on $\Ran^{\ord}(C)$.
\end{proof}

\begin{corollary}[P1.1$'$ Residue-scalar identification, lane-restricted]\label{cor:p1-1-final}
On the standard landscape (class G, L, C, M with $C_2$-cofiniteness,
non-critical, non-logarithmic), the scalar $\kappa(\mathcal A)$ is
obtained from $r(z) \in (B^{\ord})_2$ by a family-dependent
``effective residue'' prescription per Proposition~\ref{prop:kappa-dict}.
\emph{Uniform statement:} $\kappa$ is the image of $r(z)$ under
the scalar part of $\av$ at level $2$, computed via
$\av(r(z)) + \delta_{\mathrm{Sugawara}}$.
\end{corollary}

\begin{corollary}[P1.2$'$ Anyonic Yangian, lane-restricted]\label{cor:p1-2-final}
For $\mathcal A = Y(\mathfrak g)$ the Yangian, the composite
$B_n \to \Sigma_n \to 1$ has a nontrivial kernel (pure braid $P_n$);
the ordered bar $B^{\ord}$ carries a $B_n$-representation, and
$\av$ projects onto the $\Sigma_n$-invariant. The difference
$\ker(B_n \to \Sigma_n) = P_n$ encodes anyonic statistics realised
by the $R$-matrix $R(z) = 1 + \hbar\Omega/z + O(\hbar^2)$; the
``formality failure of $E_1$'' is the non-triviality of $P_n$ action
on the bar complex, which persists after $\Sigma_n$-coinvariance.
\end{corollary}

\begin{scholium}[Boltzmann-vs-Bose, pedagogical scope]\label{sch:boltzmann-pedagogy}
In the class-G lane with $\mathcal A = \mathcal H_1$ (rank-one Heisenberg),
the bar complex $B^{\ord}(\mathcal H_1)$ computes, at the Euler
characteristic level, the free-boson partition function
$\eta(\tau)^{-1}$. The ordered presentation organises by ordered mode
monomials; the symmetric presentation organises by symmetric monomials
in the oscillators. Under the Euler identity these are two
presentations of the same generating function. The mnemonic
``ordered = Boltzmann, symmetric = Bose--Einstein'' is valid in this
sense: both distribute particles into modes, and the ordered vs
symmetric choice is the choice of counting convention. Outside class
G, the mnemonic is a visual aid, not a theorem.
\end{scholium}

\section*{Summary --- The atomic theorem after five cycles}

\begin{proposition}[The atom, stated as one theorem with four scopes]\label{prop:atom-final}
The \emph{atomic phenomenon} of the chiral-bar-cobar programme is the
$\Sigma_n$-descent of Theorem~\ref{thm:p1-final}:
\[
\av: B^{\ord}(\mathcal A) \longrightarrow B^{\Sigma}(\mathcal A),
\]
realised in $\mathrm{Fact}^{\mathrm{D-mod}}(\Ran(C))$ as the
homotopy-coinvariant quotient along the $\Sigma_n$-action on the ordered
Ran-cover. Three specialisations are lane-restricted theorems:
\begin{itemize}[leftmargin=1.4em]
\item \textbf{Character identity, class G.}
$\chi(B^{\ord}) = \chi(B^{\Sigma})$ via the Euler identity on
free-boson partition functions (Theorem~\ref{thm:p1-gaussian}).
\item \textbf{Scalar extraction, class G/L/C/M.}
$\kappa = \av(r(z)) + \delta_{\mathrm{Sugawara}}$
(Proposition~\ref{prop:kappa-dict}, Corollary~\ref{cor:p1-1-final}).
\item \textbf{Anyonic braid, Yangian and $E_1$-chiral.}
$\ker(B_n \to \Sigma_n) = P_n$ acts nontrivially on $B^{\ord}$;
the $R$-matrix is the generating datum of this action
(Corollary~\ref{cor:p1-2-final}).
\end{itemize}
The Gibbs factor $1/n!$ is the universal denominator of the
$\Sigma_n$-averaging projector $\frac{1}{n!}\sum_\sigma \sigma^*$,
equivalently the normaliser of the $\Sigma_n$-cycle-index polynomial
$Z_n$, equivalently the coefficient of $x^n$ in $\exp(x)$. It is
\emph{not} a dimensional ratio between $B^{\ord}_n$ and $B^{\Sigma}_n$.
\end{proposition}

\begin{remark}[What P1 gives up and what it keeps]\label{rem:p1-concessions}
After five attack-heal cycles, the programme concedes:
\begin{enumerate}[leftmargin=1.4em]
\item ``$n!/1$ ratio'' is not the dimensional ratio; it is the Euler
denominator of averaging.
\item ``$\av$ = naive plethystic average'' is false on the vertex
algebra locus; the truth is $R$-twisted average
(concordance.tex:3980--3987).
\item ``$\av(r(z)) = \kappa$'' is family-dependent, not a uniform
residue; it requires family-specific Sugawara corrections.
\item ``$E_1$-formality failure $\equiv$ statistical gap'' is a metaphor,
sharpened to: Fresse--Willwacher non-formality is the nontriviality
of the $\hbar \ge 1$ coefficients of $R_\sigma$ in the $R$-twisted
average.
\item ``Boltzmann vs Bose-Einstein'' is correct on class G
(free-boson / lattice / Niemeier) as a presentation equivalence;
it is a pedagogical scholium, not a uniform theorem.
\end{enumerate}
And keeps:
\begin{enumerate}[leftmargin=1.4em]
\setcounter{enumi}{5}
\item The $\Sigma_n$-descent $\av: B^{\ord} \to B^{\Sigma}$ is the
atom: a symmetric-monoidal $(\infty, 1)$-functor in the factorisation
D-module category on $\Ran(C)$.
\item The image $B^{\Sigma}$ is the programme's modular object; the
five theorems A--D+H are statements about $B^{\Sigma}$ that survive
after lane-appropriate corrections.
\item The kernel of $\av$ carries the $R$-matrix, KZ associator, and
higher Yangian coherences; those are the $E_1$-specific data that
$\av$ projects away.
\item The anyonic reading on the Yangian is a separate theorem, not a
specialisation of P1; it holds on $Y(\mathfrak g)$ where the
$B_n$-representation genuinely does not factor through $\Sigma_n$.
\end{enumerate}
This five-cycle sharpening converts one overreach into four
lane-restricted survivors plus one atomic descent statement. Each
lane-restricted statement is verifiable against primary sources
(Macdonald plethystic identity, Fresse--Willwacher non-formality,
Freudenthal strange formula, Zamolodchikov conformal anomaly,
Drinfeld--Kohno theorem). The atomic descent is the
$(\infty,1)$-descent of Beilinson--Drinfeld on the factorisation
Ran-structure.
\end{remark}

\begin{remark}[Primary verifications]\label{rem:primary-verify}
Each numerical claim in Proposition~\ref{prop:kappa-dict} and
Theorem~\ref{thm:p1-gaussian} is cross-checked against the master
table in landscape\_census.tex:
\[
\kappa(\mathcal H_k) = k,\quad
\kappa(V_k(\mathfrak g)) = \frac{(k+h^\vee)\dim(\mathfrak g)}{2h^\vee},\quad
\kappa(\mathrm{Vir}_c) = c/2,\quad
\kappa(\mathcal W_N) = c(H_N - 1).
\]
$\dim V^{\otimes n} = d^n$ and $\dim \Sym^n V = \binom{d+n-1}{n}$ are
verified by direct combinatorial counting. The Macdonald-Liao identity
$\sum_n q^n \dim \Sym^n V = \prod_k (1 - q^k)^{-d}$ is the infinite
product presentation of the free-boson character; the ratio
$\binom{d+n-1}{n} / d^n \to 1/n!$ (not $n! \to 1/n!$!) as $n$ fixed,
$d \to \infty$ is the bare Stirling asymptotics, and records the
Gibbs normaliser as a limiting coefficient, not a global ratio.
\end{remark}

\begin{remark}[What was wrong in the original P1]\label{rem:original-overreach}
The original P1 \emph{overreach} is of type: stating one theorem where
the underlying mathematical content is four theorems at four different
scopes. The overreach propagates through the programme manifesto by
asserting that ``the atom'' is physically interpretable as
Boltzmann-to-Bose. After Cycle 5, the correct statement is:
\begin{itemize}[leftmargin=1.4em]
\item The atom is the $(\infty,1)$-categorical $\Sigma_n$-descent
(Theorem~\ref{thm:p1-final}).
\item The physical reading is a corollary on class G, using the
Euler identity.
\item The scalar extraction is a family-specific reading.
\item The anyonic reading is a Yangian-specific reading.
\end{itemize}
The programme benefits from the attack-heal spiral: the atom becomes
smaller (just a descent), three corollaries become sharper
(lane-restricted), and the overreach is replaced by honest scope
declaration. Beilinson's dictum applies: smaller true statements
beat a larger false one.
\end{remark}

\section*{Appendix. Cross-checks against landscape\_census.tex}

Every numerical value in this document is reconciled with
\texttt{chapters/examples/landscape\_census.tex}:
\begin{itemize}[leftmargin=1.4em]
\item Heisenberg $\mathcal H_k$: $c = 1$, $\kappa = k$, class G,
$r_{\max} = 2$; master table row 137, shadow table row 513.
\item Affine $V_k(\mathfrak g)$ general: $c = kd/(k+h^\vee)$,
$c + c' = 2d$, $\kappa = (k+h^\vee)d/(2h^\vee)$; row 146.
\item $\widehat{\mathfrak{sl}}_2$: $c = 3k/(k+2)$, $\kappa = 3(k+2)/4$;
row 150.
\item Virasoro $\mathrm{Vir}_c$: $c = 1 - 6(k+1)^2/(k+2)$, $c + c' = 26$,
$\kappa = c/2$; row 195.
\item $\mathcal W_N^k(\mathfrak{sl}_N)$: $\kappa = c\sum_{j=2}^N 1/j =
c(H_N - 1)$; row 204.
\item Shadow coefficients: $S_2 = c/2$, $S_3 = 2$, $S_4 =
10/[c(5c+22)]$, $S_5 = -48/[c^2(5c+22)]$ (proposition at line 461).
\item $\av$ mapping: concordance.tex:4026--4049 (averaging map and
$\kappa = \av(r(z)) + \dim(\mathfrak g)/2$ for Kac--Moody).
\item $R$-twisted coinvariance: concordance.tex:3980--3987
(\texttt{prop:r-matrix-descent-vol1}).
\end{itemize}
No value in this document conflicts with landscape\_census.tex or
concordance.tex. The sharpening of P1 is a \emph{scope restriction}
(Beilinson, ``prefer smaller true to larger false''), not a new
numerical assertion.

\end{document}
